\documentclass[10pt,a4paper]{amsart}
\usepackage[margin=19mm]{geometry}
\usepackage{amsmath,amssymb,mathtools}
\usepackage[scr=rsfs,cal=euler]{mathalfa}
\usepackage{xfrac}
\usepackage[unicode,hidelinks,bookmarks=false]{hyperref}
\newcommand{\ie}{i.e.\ }
\newcommand{\cf}{cf.\ }
\newcommand{\resp}{respectively\ }
\newcommand{\Subsection}{\subsection}
\providecommand{\nobreakdash}{\nobreak\hbox{-}}
\setlength{\emergencystretch}{2em}
\multlinegap0pt
\usepackage{bm}
\usepackage{bbm}
\usepackage{dsfont}
\usepackage{xspace}
\usepackage{cancel}
\usepackage{mathtools}
\usepackage{amsthm}
\makeatletter
\@ifundefined{theorem}{\newtheorem{theorem}{Theorem}}{}
\@ifundefined{lemma}{\newtheorem{lemma}[theorem]{Lemma}}{}
\makeatother
\def\IC{{\mathbb C}}
\def\IF{{\mathbb F}}
\def\IR{{\mathbb R}}
\def\bV{{\mathbf V}}
\def\span{{\mathop{\mathrm{Span}}\nolimits}}
\title[(Real)linear preservers of multiples of unitaries and matrix]{(Real)linear preservers of multiples of unitaries and matrix pairs with some extremal norm properties}
\author{Bojan Kuzma, Chi-Kwong Li, Edward Poon}
\date{Source: arXiv:2510.05303v1 (CC BY 4.0)}
\begin{document}
\maketitle

\noindent\emph{Source and licence.} (Real)linear preservers of multiples of unitaries and matrix pairs with some extremal norm properties. Version arXiv:2510.05303v1 (CC BY 4.0), distributed under the Creative Commons Attribution 4.0 International licence, as recorded in the arXiv licence metadata for this exact version and shown on the abstract page https://arxiv.org/abs/2510.05303v1. Full rights notice: licenses/arxiv-2510.05303-CC-BY-4.0.txt.\par\smallskip

\noindent\textit{Editorial scope.} Counted unit: the Lemma whose source label is \texttt{lem:aux2-by-2-reallinear} in the article (source0.tex lines 769--776, source label \texttt{lem:aux2-by-2-reallinear}), delivered together with the article's own complete original proof of it (source0.tex lines 778--834, one \texttt{proof} environment of 57 lines). No mathematics is written, repaired or re-derived: statement and proof are the article's own text line for line. Required context retained uncounted: lam1 (lines 307--313); P:unitary (lines 407--418). Every definition, display and label of those lines is retained unchanged. Omitted from this package and not redistributed: 1--306, 314--406, 419--768, 777--777, 835--1046. Nothing inside a retained excerpt is omitted or altered: every retained body is delivered exactly as the article's lines are written, so no omission receipt of any kind accompanies this package. Citations: the retained excerpts carry no citation command at all, so no bibliography entry of the article is transcribed and no reference is redistributed. Reference adaptation: none was needed and none was made; every reference inside the delivered unit resolves inside it, so no unresolved reference or citation is produced. Printed slips kept literal: no printed word, symbol, hypothesis or number of the counted statement or of its proof was altered. Layout and class: the article's own class is \texttt{article} with packages including amsmath, amssymb, amsthm, verbatim, color, graphics, srcltx, hyperref; the wrapper is the lane's pinned \texttt{amsart} editorial wrapper at 10pt on A4, which restates the theorem-like environments the retained text needs and adds the article's own macro definitions that the retained text needs (\texttt{\textbackslash IC}, \texttt{\textbackslash IF}, \texttt{\textbackslash IR}, \texttt{\textbackslash bV}, \texttt{\textbackslash span}), with extra wrapper packages (bm, bbm, dsfont, xspace, cancel, mathtools). The delivered numbering is the unit's own. Assets: no retained span contains a figure environment, a graphics-inclusion command, a PDF-page inclusion, a table or a TikZ picture; no image file is copied into this package, the package declares no asset dependency and no image file is redistributed with it. Licence: version arXiv:2510.05303v1 of this article is distributed under the Creative Commons Attribution 4.0 International licence, as recorded in the arXiv licence metadata for this exact version and shown on the abstract page https://arxiv.org/abs/2510.05303v1. A full rights notice is delivered at licenses/arxiv-2510.05303-CC-BY-4.0.txt. Characters: every retained excerpt is pure ASCII, so the delivered engine, XeLaTeX with its default Latin Modern text font, reports no missing character.
 \begin{lemma}\label{lam1} The following is equivalent for a matrix $A\in M_n(\IF)$.
 \begin{itemize}
  \item[{\rm (i)}] $\|AB\|=\|A\|\cdot\|B\|$ for every $B\in M_n(\IF)$
   \item[{\rm (ii)}] $\|AB^\ast\|=\|A\|\cdot\|B^\ast\|$ for every $B\in M_n(\IF)$
   \item[{\rm (iii)}] $\gamma A\in {\mathcal U}_n(\IF)$ for some $\gamma \in \IF$.
 \end{itemize}
 \end{lemma}

\begin{theorem}\label{P:unitary}
A bijective real-linear map $T\colon  M_2(\IF) \rightarrow M_2(\IF)$
maps matrices in ${\mathcal U}_2(\IF)$ to multiples of matrices in ${\mathcal U}_2(\IF)$ if and only if $T(I)$ is a nonzero multiple of a  matrix  in ${\mathcal U}_2(\IF)$ and the map $L$ defined by $L(X) = T(I)^{-1} T(X)$ satisfies one of the following.
\begin{itemize}
\item[{\rm (i)}] $\IF = \IR$ and $L(\bV_j) = \bV_j$ for $j = 1,2$, where
$\bV_1 = \span\{I_2, E_{12}-E_{21}\}$ and $\bV_2 = \span\{E_{11}-E_{22}, E_{12} + E_{21}\}$.
\item[{\rm (ii)}] $\IF=\IC$ and there are $U,V\in {\mathcal U}_2(\IC)$ and $\mu\in\IC\setminus\IR$ such that the unital map $L'\colon X \mapsto
 VL(UXU^\ast)V^\ast$ satisfies $L'(i I)=\mu I$ and
               $$L'(\Sigma_j) = a_j \Sigma_j +  b_j I_2\quad \hbox{and}\quad L'(i\Sigma_j)=\mu L'(\Sigma_j)$$     for some
real $a_1, a_2, a_3, b_1, b_2, b_3$ with  $a_1\ge a_2\ge| a_3|>0$. If $T$ is complex-linear, then $\mu=i$.
\end{itemize}
\end{theorem}

\begin{lemma}\label{lem:aux2-by-2-reallinear}
 	Suppose $T\colon M_2(\IC) \rightarrow M_2(\IC)$ is a  unital real-linear bijection which satisfies any one among the identities
    \begin{itemize}
     \item[{\rm (i)}] $\|T(A)T(B)\| = \|T(A)\| \|T(B)\| \qquad \hbox{ whenever } \qquad \|AB\| = \|A\| \|B\|$.
     \item[{\rm (ii)}] $\|T(A)T(B)^\ast\| = \|T(A)\| \|T(B)^\ast\| \qquad \hbox{ whenever } \qquad \|AB^\ast\| = \|A\| \|B^\ast\|$.
      \end{itemize}
	Then there exist  unitary $U,V$ such that $L\colon X\mapsto  V T(U XU^\ast)V^\ast$ fixes all diagonal matrices in $\bV_3$ and multiplies off-diagonal ones in $\bV_3$  by some $w>0$.
\end{lemma}

\begin{proof}
    By Lemma~\ref{lam1}, $T$ maps unitary matrices into scalar multiples of unitary matrices.
Using the unital assumption, we may apply Theorem~\ref{P:unitary} (ii)
 and modify $T$ to the map $X \mapsto VT(UXU^*)V^*$ with $U, V \in U_2(\IC)$ and assume that

\begin{equation}\label{eq:AB*-on-2-by-2anew}
\begin{aligned}
  &T \left( \left[\begin{smallmatrix} i & 0 \\ 0 & -i \end{smallmatrix}\right] \right) = a \left[\begin{smallmatrix} i & 0 \\ 0 & -i \end{smallmatrix}\right] +bI,\quad T \left( \left[\begin{smallmatrix} 0 & i \\ i & 0 \end{smallmatrix}\right] \right) = w\left[\begin{smallmatrix} 0 & i \\ i & 0 \end{smallmatrix}\right] +cI, \quad T \left( \left[\begin{smallmatrix} 0 & 1 \\ -1 & 0 \end{smallmatrix}\right] \right) = z\left[\begin{smallmatrix} 0 & 1 \\ -1 & 0 \end{smallmatrix}\right] + dI\\
  &T(iX)=\mu T(X);\qquad X\in\bV_3
\end{aligned}
\end{equation}
for some $\mu\in\IC\setminus\IR$ and some real constants satisfying $a\ge w\ge|z| > 0$. Notice that, due to bijectivity and unitality, $a, w, z$  are nonzero, so $a,w>0$.


Let $A_r = \begin{bmatrix} (r+1)i& 0 \\ 0 & (r-1)i \end{bmatrix}$, $r>0$. By~\eqref{eq:AB*-on-2-by-2anew} and as $T$ is unital,
$$T(A_r)=(b+r\mu) I +a\begin{bmatrix} i &0 \\ 0 & -i \end{bmatrix} = \begin{bmatrix}
 i a +b +r\mu  & 0 \\
 0 & -i a +b +r\mu
 \end{bmatrix} $$
Write $\mu=|\mu|e^{i\theta}$ and assume first that $\sin\theta>0$. Since $a,r>0$ and $b\in\IR$, one computes that
$A_r=-A_r^\ast$, $T(A_r)$, and $T(A_r)^\ast$ all
 attain their norms only at scalar multiples of their eigenvector $e_1$ (their $(11)$-entry has larger modulus than their $(22)$-entry).  Notice that $\|XA_r\| = \|X \| \|A_r\|$  for all $r > 0$ if and only if $X$
attains its norm on the vector $e_1$.
 This is equivalent to its SVD equaling $X =x e_1^\ast+ye_2^\ast$ for some orthogonal $x,y$ with $\|x\|\ge \|y\|$,  that is, $X$ must have the form
\begin{equation}\label{e1norm1anew}
X = \begin{bmatrix} \alpha & t \bar{\beta} \\ \beta & -t\bar{\alpha} \end{bmatrix}
\end{equation}
for some $\alpha, \beta, t \in \IC$ with $|t| \leq 1$. Thus, assuming either (i) or (ii) in the statement of this lemma, $X$ from~\eqref{e1norm1anew} is mapped into $T(X)$ which must also attain its norm at $e_1$ and hence have the same form~\eqref{e1norm1anew}.

However if, with $\mu=|\mu|e^{i\theta}$, one has $\sin\theta<0$, then  one computes that $T(A_r)$ and $T(A_r)^*$ both attain their
norm only at scalar multiples of their eigenvector $e_2$. Proceeding as above one sees that, under either of the assumptions (i) or (ii), $X$ of the form~\eqref{e1norm1anew} is mapped into a matrix $T(X)$ which attains its norm at $e_2$, hence is of the same form  as in~\eqref{e1norm1anew}  except that $|t| >1$. Hence, regardless if $\sin\theta$ is positive or negative, $X$ of the form~\eqref{e1norm1anew} is mapped into a matrix $T(X)$ with  its first  column orthogonal to its  second one.
 If we set $\alpha=\cos\pi/4$ and $\beta = e^{i\theta}\sin\pi/4$ and $t=1/3$, then under assumptions (i)--(ii),
$$T(X) = \frac{1}{6 \sqrt{2}}
\begin{bmatrix}
 x & y \\
 u & v \\
\end{bmatrix}$$
where
$$x=-4 i a \mu -4 b \mu -(c+i d) e^{-i \theta } (2 \mu -i)-(c-i d) e^{i \theta } (2 \mu +i)+2,$$
$$y =e^{i \theta } (w-z) (1-2 i \mu
   )-i e^{-i \theta } (w+z) (2 \mu -i),$$
$$u=e^{i \theta } (w+z) (1-2 i \mu )+e^{-i \theta } (z - w) (1 + 2 i  \mu),$$
$$v=4 i a \mu -4 b \mu -(c+i d) e^{-i \theta } (2 \mu
   -i)-(c-i d) e^{i \theta } (2 \mu +i)+2.$$

Since  the columns of $T(X)$  must be orthogonal, so
\begin{equation}\label{orth-cols1anew}
	x \bar{y} + u \bar{v} = 0.
\end{equation}
The left-hand side of the above equation is a trigonometric polynomial in $e^{i\theta}$:
 $$-8(\mu-\mathrm{Re}\,\mu  )  \left(i e^{-i \theta } (a-i b+1) (w-z)-e^{i \theta } (b+i (a-1)) (w+z)-2 (c z+i d w)\right)=0.$$
Its coefficients therefore vanish identically and since, by the assumptions $\mu\notin\IR$, $a,w>0$, and all other variables are real with $z \ne 0$, the only possibility is
$$a=1,\quad b=0,\quad z=w>0,\quad d=c=0.$$
This shows that $T$ is the identity on diagonal matrices from $\bV_3$, while on off-diagonal matrices from $\bV_3$ it acts as
$$T(\Sigma_1)=w \Sigma_1,\quad T(\Sigma_2)=w\Sigma_2,$$
as claimed.
\end{proof}

\end{document}
